\begin{examplebox}

	\textbf{\textcolor{CExample}{例 1（基础）}}

	\textbf{\textcolor{CExample}{题目：}}
	已知 $a,b \in \mathbb{R}$，求证 $a^2+b^2 \ge 2ab$。

	\textbf{\textcolor{CExample}{解题步骤：}}
	\begin{description}[style=nextline, leftmargin=2.8em, labelsep=0.5em, itemsep=0.45em]
		\item[\textbf{第一步（作差变形）：}]
			\[
				a^2+b^2-2ab = (a-b)^2.
			\]
			\par\textbf{理由：}
			完全平方公式。

		\item[\textbf{第二步（判断符号）：}]
			\[
				(a-b)^2 \ge 0.
			\]
			\par\textbf{理由：}
			实数的平方为非负数，故 $(a-b)^2 \ge 0$ 对所有实数成立。

		\item[\textbf{第三步（得出结论）：}]
			\[
				a^2+b^2-2ab \ge 0 \Longrightarrow a^2+b^2 \ge 2ab.
			\]
			\par\textbf{理由：}
			由差非负（$a^2+b^2-2ab \ge 0$）根据比较法定义即得 $a^2+b^2 \ge 2ab$。
	\end{description}

	\textbf{\textcolor{CExample}{关键结论：}}
	作差后配方，利用平方非负性得证，体现比较法直接有效的特点。

	\medskip
	\textbf{\textcolor{CExample}{例 2（稍变形）}}

	\textbf{\textcolor{CExample}{题目：}}
	已知 $a,b \ge 0$，求证 $\sqrt{ab} \le \frac{a+b}{2}$。

	\textbf{\textcolor{CExample}{解题步骤：}}
	\begin{description}[style=nextline, leftmargin=2.8em, labelsep=0.5em, itemsep=0.45em]
		\item[\textbf{第一步（分析法转化）：}]
			要证 $\sqrt{ab} \le \frac{a+b}{2}$，即证 $2\sqrt{ab} \le a+b$。
			\[
				\sqrt{ab} \le \frac{a+b}{2} \;\Longleftrightarrow\; 2\sqrt{ab} \le a+b.
			\]
			\par\textbf{理由：}
			两边乘以2，不等号方向不变。

		\item[\textbf{第二步（平方去根号）：}]
			即证 $(\sqrt{a}-\sqrt{b})^2 \ge 0$。
			\[
				\begin{aligned}
					2\sqrt{ab} \le a+b &\iff a+b-2\sqrt{ab} \ge 0 \\
					&\iff (\sqrt{a}-\sqrt{b})^2 \ge 0.
			\end{aligned}
			\]
			\par\textbf{理由：}
			移项后左边是完全平方。

		\item[\textbf{第三步（结论肯定）：}]
			因为 $(\sqrt{a}-\sqrt{b})^2 \ge 0$ 恒成立，所以原不等式成立。
			\par\textbf{理由：}
			分析法由结论逆向追溯至显然成立的命题，从而确认原命题真。
	\end{description}

	\textbf{\textcolor{CExample}{关键结论：}}
	分析法通过寻找充分条件简化不等式，最后归结为平方非负。

	\medskip
	\textbf{\textcolor{CExample}{例 3（综合应用）}}

	\textbf{\textcolor{CExample}{题目：}}
	已知 $a,b \in \mathbb{R}$，$a+b > 2$，求证 $a$ 和 $b$ 中至少有一个大于 $1$。

	\textbf{\textcolor{CExample}{解题步骤：}}
	\begin{description}[style=nextline, leftmargin=2.8em, labelsep=0.5em, itemsep=0.45em]
		\item[\textbf{第一步（假设反面）：}]
			假设 $a\le 1$ 且 $b\le 1$。
			\[
				\neg(a>1 \lor b>1) \Longleftrightarrow (a\le 1 \land b\le 1).
			\]
			\par\textbf{理由：}
			反证法需先否定结论，即“至少有一个大于1”的反面是“两个都不大于1”。

		\item[\textbf{第二步（推出矛盾）：}]
			由 $a\le1$ 且 $b\le1$ 得 $a+b \le 2$。
			\[
				a\le1,\; b\le1 \;\Longrightarrow\; a+b \le 2.
			\]
			\par\textbf{理由：}
			同向不等式相加：$a\le1$ 与 $b\le1$ 相加得 $a+b\le2$。

		\item[\textbf{第三步（矛盾形成）：}]
			已知 $a+b > 2$，而推得 $a+b \le 2$，矛盾。
			\par\textbf{理由：}
			假设导致与已知条件冲突，故假设不成立，原结论成立（$a$ 和 $b$ 中至少有一个大于1）。
	\end{description}

	\textbf{\textcolor{CExample}{关键结论：}}
	反证法通过假设结论不成立，结合已知条件推出矛盾，从而证明原命题正确。
\end{examplebox}
